% !TEX program = lualatex
% Compile from 2026/week3: latexmk -lualatex slide.tex
% Source: ../../2025/week3/week3.pdf (13 slides, 4:3).
\documentclass[12pt,aspectratio=43]{beamer}
% Shared Beamer design. Load after \documentclass using \input.
% Engine: LuaLaTeX. Fonts: UD Digi Kyokasho N (regular/bold), Arial.
\usepackage{luatexja-fontspec}
\setmainjfont[BoldFont={UD Digi Kyokasho N Bold}]{UD Digi Kyokasho N}
\setsansjfont[BoldFont={UD Digi Kyokasho N Bold}]{UD Digi Kyokasho N}
\setmonojfont[BoldFont={UD Digi Kyokasho N Bold}]{UD Digi Kyokasho N}
% Liberation Sans keeps the layout usable on systems without Arial.
\IfFontExistsTF{Arial}{
  \setmainfont{Arial}
  \setsansfont{Arial}
  \setmonofont{Arial}
}{
  \setmainfont{Liberation Sans}
  \setsansfont{Liberation Sans}
  \setmonofont{Liberation Sans}
}
\renewcommand{\kanjifamilydefault}{\gtdefault}

% Typeset mathematics with a dedicated TeX math font, independent of Arial.
\usefonttheme{professionalfonts}
\usepackage[mathrm=sym]{unicode-math}
\setmathfont{Latin Modern Math}
\setoperatorfont{\symup}

\definecolor{courseblue}{HTML}{4472C4}
\definecolor{ruleblue}{HTML}{8FAADC}
\setbeamercolor{normal text}{fg=black,bg=white}
\setbeamercolor{structure}{fg=courseblue}
\setbeamercolor{frametitle}{fg=black,bg=white}
\setbeamercolor{title banner}{fg=white,bg=courseblue}
\setbeamerfont{frametitle}{size=\fontsize{18}{22}\selectfont,series=\mdseries}
\setbeamerfont{quote}{shape=\upshape}
\setbeamersize{text margin left=9mm,text margin right=9mm}
\setbeamertemplate{navigation symbols}{}
\setbeamercolor{page number}{fg=gray,bg=white}
\setbeamertemplate{footline}{%
  \begin{beamercolorbox}[wd=\paperwidth,ht=2.5mm,dp=1.5mm]{page number}
    \hfill{\fontsize{8}{9}\selectfont
      \ifnum\value{framenumber}>0
        \insertframenumber\,/\,\inserttotalframenumber
      \fi}\hspace*{3mm}
  \end{beamercolorbox}}
\setbeamertemplate{headline}{%
  \hfill\textcolor{gray}{\fontsize{6}{7}\selectfont
    \hypersetup{linkcolor=gray}% Keep Beamer's title links in the header color.
    \ifnum\value{framenumber}=0 \insertshortdate\enspace\fi
    \insertshorttitle\space\insertshortsubtitle}\hspace{2mm}\vspace{1.5mm}}
\setbeamertemplate{frametitle}{%
  \vspace{2mm}\centering\insertframetitle\par
  \vspace{-1mm}{\color{ruleblue}\rule{.44\paperwidth}{.5pt}}\par
  \vspace{2mm}}
\setbeamercolor{itemize item}{fg=courseblue}
\setbeamercolor{itemize subitem}{fg=black}
\setbeamercolor{itemize subsubitem}{fg=black}
\setbeamertemplate{itemize item}{\textbullet}
\setbeamertemplate{itemize subitem}{\textnormal{-}}
\setbeamertemplate{itemize subsubitem}{\textnormal{-}}
\setlength{\parskip}{.45em}
\hypersetup{colorlinks=true,urlcolor=courseblue}

% Use \titlepage inside a [t,noframenumbering] frame. The background extends
% 1 mm past both page edges to avoid hairline gaps in PDF viewers.
\setbeamertemplate{title page}{%
  \vspace{2mm}
  \begin{beamercolorbox}[wd=\paperwidth,ht=43mm,dp=13mm,colsep*=1mm,center]{title banner}
    {\fontsize{32}{40}\selectfont\inserttitle\par}
    \vspace{2mm}
    {\fontsize{20}{24}\selectfont\insertsubtitle\par}
  \end{beamercolorbox}
  \vspace{10mm}
  \hfill\insertauthor
}

\usepackage{listings}
\usepackage{upquote}
\IfFontExistsTF{Consolas}{\newfontfamily\codefont{Consolas}}{\newfontfamily\codefont{Liberation Mono}}
\definecolor{codevalue}{HTML}{7A3E9D}
\definecolor{codebuiltin}{HTML}{006DB0}
\definecolor{codestring}{HTML}{267348}
\newcommand{\code}[1]{{\codefont #1}}

% Python examples remain editable, with spaces and punctuation preserved.
\lstdefinestyle{python}{
  language=Python,
  basicstyle=\codefont\fontsize{14}{16}\selectfont,
  keywordstyle=\color{codebuiltin},
  keywordstyle=[2]\color{codebuiltin},
  morekeywords=[2]{print,range},
  stringstyle=\color{codestring},
  showstringspaces=false,
  columns=fullflexible,
  keepspaces=true,
  upquote=true,
  tabsize=4,
  breaklines=false,
  backgroundcolor=\color{black!4},
  frame=single,framerule=0pt,framesep=1.7mm,
  xleftmargin=1.7mm,xrightmargin=1.7mm,
  aboveskip=1mm,belowskip=0.5mm,
  literate=*{0}{{{\color{codevalue}0}}}1
    {1}{{{\color{codevalue}1}}}1 {2}{{{\color{codevalue}2}}}1
    {3}{{{\color{codevalue}3}}}1 {4}{{{\color{codevalue}4}}}1
    {5}{{{\color{codevalue}5}}}1 {6}{{{\color{codevalue}6}}}1
    {7}{{{\color{codevalue}7}}}1 {8}{{{\color{codevalue}8}}}1
    {9}{{{\color{codevalue}9}}}1
}
\lstset{style=python}
\newcommand{\result}[1]{%
  \par\vspace{1mm}%
  {\small\textcolor{black!60}{出力}}\quad
  {\codefont\fontsize{16}{18}\selectfont #1}\par}
\setbeamercolor{lesson note}{fg=black,bg=courseblue!9}
\newenvironment{lessonbox}
  {\begin{beamercolorbox}[wd=\linewidth,sep=1.5mm]{lesson note}}
  {\end{beamercolorbox}\par}

\title{計算物理学B}
\subtitle{第三回}
\author{野垣 康介、藤本 悠輝}
\date{2026/11/06}
\hypersetup{pdftitle={計算物理学B 第三回},pdfauthor={野垣 康介、藤本 悠輝}}

\begin{document}

% 1. 表紙
\begin{frame}[t,noframenumbering]
  \titlepage
\end{frame}

% 2. 2026年度の講義予定。11/6（金）の関数を強調。
\begin{frame}{講義予定}
  \setlength{\parskip}{0pt}
  2026年度 第2学期（火曜3限）
  \par\vspace{2mm}
  \begingroup\centering
    \renewcommand{\arraystretch}{1.05}
    \begin{tabular}{@{}l@{\hspace{2mm}}l@{\hspace{4mm}}l@{\hspace{2mm}}l@{}}
      10/6 & 四則演算 & 12/8 & モンテカルロ1 \\
      10/27 & 制御文（for, if） & 12/15 & モンテカルロ2 \\
      \textcolor{courseblue}{\textbf{11/6（金）}} &
      \textcolor{courseblue}{\textbf{関数}} & 12/22 & 微分方程式1 \\
      11/10 & 配列（numpy） & 1/12 & 微分方程式2 \\
      11/17 & 可視化（matplotlib） & 1/19 & 偏微分方程式1 \\
      11/24 & 数値微分 & 1/26 & 偏微分方程式2 \\
      12/1 & 数値積分 & 2/2 & 根の探索 \\
      \multicolumn{2}{c}{\small 〜中間レポート〜} &
      \multicolumn{2}{c}{\small 〜期末レポート〜}
    \end{tabular}\par
  \endgroup
  \vspace{2mm}
  {\small
    \textcolor{courseblue}{10/20は出張で休講。}\\
    11/6（金）は火曜授業の振替。\par}
  \vspace{1mm}
  {\fontsize{7}{8}\selectfont
    日付：\href{https://www.niigata-u.ac.jp/wp-content/uploads/2025/08/schedule_2026.pdf}{令和8年度（2026）新潟大学授業暦}\par}
\end{frame}

% 3. 授業で使う教材
\begin{frame}{授業で使うリンク}
  \begin{itemize}
    \setlength{\itemsep}{6mm}
    \item Google Colab\\
      {\small\url{https://colab.research.google.com/?hl=ja}}
    \item 講義資料（GitHub）\\
      {\small\url{https://github.com/nogaki/Computational_Physics_B}}
    \item 今週の教材
      \begin{itemize}
        \item フォルダー：\code{2026/week3}
        \item 実習用ノートブック：\code{week3.ipynb}
      \end{itemize}
  \end{itemize}
\end{frame}

% 4. week2と同じColabへの取り込み・保存手順
\begin{frame}[t]{GitHubからColabへ取り込む}
  \setlength{\parskip}{0pt}
  \begin{enumerate}
    \setlength{\itemsep}{4mm}
    \item GitHubで\code{week3.ipynb}を開き、\\
      ブラウザーのアドレス欄のURLをコピーします。
    \item Colabで「ファイル」→「ノートブックを開く」\\
      →「GitHub」を選びます。
    \item コピーしたURLを貼り付けて検索し、\\
      表示されたノートブックを開きます。
  \end{enumerate}
  \vspace{3mm}
  \begin{lessonbox}
    \textbf{開いたら「ドライブにコピーを保存」}\\[1mm]
    \small 「ファイル」メニューから、自分用のコピーを保存します。\\
    開いただけでは、自分のドライブに保存されません。
  \end{lessonbox}
\end{frame}

% 4. 関数を使う前に、重複する処理を見る
\begin{frame}[t,fragile]{関数：同じ処理の繰り返し}
  \setlength{\parskip}{0pt}
  1から10までと、1から30までの和をそれぞれ求めます。
\begin{lstlisting}
sum1 = 0
for i in range(10):
    sum1 = sum1 + (i + 1)

sum2 = 0
for j in range(30):
    sum2 = sum2 + (j + 1)

print(sum1 + sum2)
\end{lstlisting}
  {\small ほぼ同じ処理を、範囲だけ変えて2回書いています。}
\end{frame}

% 5. 重複する総和処理を関数にする
\begin{frame}[t,fragile]{関数：処理をまとめる}
  \setlength{\parskip}{0pt}
\begin{lstlisting}
def sum_one2n(n):
    total = 0
    for i in range(n):
        total = total + (i + 1)
    return total
sum1 = sum_one2n(10)
sum2 = sum_one2n(30)
print(sum1 + sum2)
\end{lstlisting}
  \result{520}
  \vspace{1mm}
  \begin{lessonbox}
    \small 総和を求める処理を関数にまとめ、引数を変えて呼び出します。
  \end{lessonbox}
\end{frame}

% 6. 関数の定義と呼び出し
\begin{frame}[t,fragile]{関数の定義と返り値}
  \setlength{\parskip}{0pt}
\begin{lstlisting}
def function_name(argument):
    return argument
\end{lstlisting}
  \begin{itemize}
    \setlength{\itemsep}{2mm}
    \item \code{def}の後に関数名と引数を書き、行末に\code{:}を付けます。
    \item 関数の中の処理は字下げします（空白4個を目安に）。
    \item 引数は複数指定できます。\code{return}で呼び出し元に値を返します。
    \item 返り値が複数の場合は、タプル（\code{tuple}）として返ります。
    \item 関数内で定義した変数は、関数の外では使えません。
  \end{itemize}
\end{frame}

% 7. printとreturnの違い
\begin{frame}[t,fragile]{printとreturnの違い}
  \setlength{\parskip}{0pt}
  \code{print}は表示し、\code{return}は値を返します。\code{return}がない関数の返り値は\code{None}（値なし）です。
  \begin{columns}[T,onlytextwidth]
    \begin{column}{.48\textwidth}
      \textbf{\code{print}で表示}
\begin{lstlisting}[aboveskip=0mm,belowskip=0mm]
def show(x):
    print(x * x)

a = show(3)
print(a)
\end{lstlisting}
      {\small\textcolor{black!60}{出力}}\quad
      {\codefont\fontsize{16}{18}\selectfont 9\\
        \hspace*{1cm}None\par}
      {\small \code{a}には\code{None}が入ります。}
    \end{column}
    \begin{column}{.48\textwidth}
      \textbf{\code{return}で値を返す}
\begin{lstlisting}[aboveskip=0mm,belowskip=0mm]
def square(x):
    return x * x

b = square(3)
print(b)
\end{lstlisting}
      \result{9}
      {\small \code{b}に入る\code{9}は後の計算に使えます。}
    \end{column}
  \end{columns}
\end{frame}

% 8. ローカル変数tmpを関数の外から参照するとNameError
\begin{frame}[t,fragile]{関数内の変数}
  \setlength{\parskip}{0pt}
\begin{lstlisting}
def sum_all(x, y, z):
    tmp = x + y
    return tmp + z

l = 3
m = 2
n = 5
print(sum_all(l, m, n))
\end{lstlisting}
  \result{10}
  \vspace{1mm}
  \small 関数内で作った\code{tmp}は外から使えません。\code{print(tmp)}を実行すると、意図的に\code{NameError}になります。
\end{frame}

% 9. 引数と物理パラメータ
\begin{frame}[t,fragile]{引数と物理パラメータ}
  \setlength{\parskip}{0pt}
  運動エネルギーは\(E=\frac{1}{2}mv^2\)です。\(m\)は質量、\(v\)は速度です。
  \par\vspace{1mm}
  単位は、質量\(m\)：\(\mathrm{kg}\)、速度\(v\)：\(\mathrm{m/s}\)、エネルギー\(E\)：\(\mathrm{J}\)です。
\begin{lstlisting}
def kinetic_energy(m, v):
    return 0.5 * m * v ** 2

print(kinetic_energy(2.0, 3.0))
\end{lstlisting}
  \result{9.0}
  \begin{lessonbox}
    \small 呼び出しの1番目の値\code{2.0}が\code{m}、2番目の値\code{3.0}が\code{v}に入ります。
  \end{lessonbox}
\end{frame}

% 10. 自由落下で複数の値を返す
\begin{frame}[t,fragile]{複数の返り値：自由落下}
  \setlength{\parskip}{0pt}
  空気抵抗を無視し、静止状態から落下する物体を考えます。落下開始位置を\(x=0\)、下向きを正とします。
  \par\vspace{1mm}
  計算例では重力加速度を\(g=10.0\ \mathrm{m/s^2}\)と近似します。経過時間\(t\geq0\)は秒、位置\(x\)はメートル、速度\(v\)は\(\mathrm{m/s}\)です。
  \[x=\frac{1}{2}gt^2,\qquad v=gt\]
\begin{lstlisting}
def free_fall(t):
    g = 10.0
    x = 0.5 * g * t ** 2
    v = g * t
    return x, v
\end{lstlisting}
\end{frame}

% 11. タプルを受け取る
\begin{frame}[t,fragile]{複数の返り値：タプルを受け取る}
  \setlength{\parskip}{0pt}
  カンマでまとめた値の組をタプルといいます。\code{return x, v}は、位置と速度をこの順でまとめて返します。
  \par\vspace{1mm}
  返した順で複数の変数に受け取ります。変数の数は返す値の数に合わせます。
\begin{lstlisting}
state = free_fall(3.0)
print(state)
x, v = state
print(x, v)
\end{lstlisting}
  {\small\textcolor{black!60}{出力}}\quad
  {\codefont\fontsize{16}{18}\selectfont
    (45.0, 30.0)\\
    \hspace*{1cm}45.0\quad 30.0\par}
  \begin{lessonbox}
    \small \code{x}は\(45.0\ \mathrm{m}\)、\code{v}は\(30.0\ \mathrm{m/s}\)です。
  \end{lessonbox}
\end{frame}

% 12. 関数を組み合わせる
\begin{frame}[t,fragile]{関数を組み合わせる}
  \setlength{\parskip}{0pt}
  \code{free\_fall}が返す速度を\code{kinetic\_energy}の引数に渡し、運動エネルギーを計算します。
\begin{lstlisting}
m = 2.0
x, v = free_fall(3.0)
energy = kinetic_energy(m, v)
print(energy)
\end{lstlisting}
  \result{900.0}
  \begin{lessonbox}
    \small 質量\(2.0\ \mathrm{kg}\)、落下開始から\(3.0\ \mathrm{s}\)後の運動エネルギーは\(900.0\ \mathrm{J}\)です。
  \end{lessonbox}
\end{frame}

% 13. 簡単な入力で結果を確かめる
\begin{frame}[t,fragile]{計算結果を確かめる}
  \setlength{\parskip}{0pt}
  実行前に、簡単な入力で結果を予想します。
\begin{lstlisting}
print(kinetic_energy(2.0, 0.0))
print(free_fall(0.0))
\end{lstlisting}
  \result{0.0}
  \result{(0.0, 0.0)}
  \small 静止していれば運動エネルギーは0です。落下開始時\(t=0\)の位置と速度も0です。
  \par\vspace{2mm}
  \textbf{短い実習}　基準を\(m=2.0\ \mathrm{kg},\ v=3.0\ \mathrm{m/s}\)とします。質量だけを2倍、速度だけを2倍にした場合のエネルギーを別々に計算し、基準の何倍か説明しましょう。
\end{frame}

% 14. 漸化式を反復計算する
\begin{frame}[t,fragile]{関数の活用：漸化式}
  \setlength{\parskip}{0pt}
  漸化式 \(a_1=1,\quad a_{n+1}=2a_n+3\) の一般項は
  \(a_n=4\cdot2^{n-1}-3\) です。\(a_{10}\)を求めます。
\begin{lstlisting}
def next_term(a):
    return 2 * a + 3

a = 1
for _ in range(9):
    a = next_term(a)

print(a)
\end{lstlisting}
  \result{2045}
\end{frame}

% 15. ユークリッドの互除法をwhile文で書く
\begin{frame}[t,fragile]{最大公約数：ループ版}
  \setlength{\parskip}{0pt}
  \begin{columns}[T,onlytextwidth]
    \begin{column}{.49\textwidth}
\begin{lstlisting}
def gcd(n, m):
    while m != 0:
        r = n % m
        n = m
        m = r
    return n
\end{lstlisting}
    \end{column}
    \begin{column}{.48\textwidth}
      \small
      ユークリッドの互除法：\\
      正の整数\(n,m\)に対し、\\
      余り\(r=n\bmod m\)を計算し、\\
      \((n,m)\)を\((m,r)\)に更新。\\
      \(m=0\)になるまで繰り返し、\\
      最後の\(n\)が最大公約数です。
    \end{column}
  \end{columns}
\begin{lstlisting}[aboveskip=1mm,belowskip=0mm]
print(gcd(120, 12))
\end{lstlisting}
  \result{12}
\end{frame}

% 16. 原資料の例題と2問。演習の答えは載せない。
\begin{frame}[t,fragile]{実習タイム}
  \setlength{\parskip}{0pt}
  \textbf{例題}　関数を使って次の値を求めましょう。
  \[\displaystyle \sum_{n=1}^{10} a_n,\qquad a_n=n^2+4n+5\]
  \vspace{-2mm}
  \begin{enumerate}
    \setlength{\itemsep}{2mm}
    \item \textbf{フィボナッチ数列}の第10項を、関数を使って求めましょう。\\
      \(F_1=F_2=1,\quad F_{n+2}=F_{n+1}+F_n\)
    \item 引数\(n\)が素数かどうか判定する関数を作りましょう。
  \end{enumerate}
  \vspace{1mm}
  \small 素数判定のヒント：\code{math.sqrt(n)}で平方根、\code{math.floor(x)}で\(\lfloor x\rfloor\)を計算できます。
\begin{lstlisting}[aboveskip=1mm,belowskip=0mm]
import math
\end{lstlisting}
\end{frame}

% 17. 発展：再帰関数
\begin{frame}[t,fragile]{発展：最大公約数（再帰版）}
  \setlength{\parskip}{0pt}
  ループ版と同じ処理を、関数が自分自身を呼び出す再帰で書けます。
\begin{lstlisting}
def gcd_recursive(n, m):
    if m == 0:
        return n
    return gcd_recursive(m, n % m)

print(gcd_recursive(120, 12))
\end{lstlisting}
  \result{12}
  \begin{lessonbox}
    \small \code{m == 0}が再帰の終了条件です。そうでなければ、次の組\code{(m, n \% m)}を計算します。
  \end{lessonbox}
\end{frame}

\end{document}
